\documentclass[12pt]{article}
\usepackage[utf8]{inputenc}
\usepackage{setspace}
\usepackage{geometry}
\usepackage{hyperref}
\usepackage{graphicx}
\usepackage{float}


\geometry{margin=1in}
\setstretch{1.25}

\title{CSP 571 Project Phase 2: Crypto Tweet Sentiment and Price Movement Analysis}
\author{Madhu Siddharth Suthagar - A20567167 \\ John Chmura - A20507144 \\ Akash Chenchugan - A20583620 \\ Harish Namasivayam Muthuswamy - A20588339 \\ Yogesh Periyasamy - A20570447}
\date{November 12, 2025}

\begin{document}

\maketitle

\section{Project Overview}
We aim to investigate whether Elon Musk's cryptocurrency-related tweets, especially those mentioning Dogecoin and Bitcoin, are associated with subsequent price movements in DOGE/USD and BTC/USD. The analysis will explore whether social media sentiment can act as a short-term or medium-term indicator for cryptocurrency returns.

\section{Data Integration Steps}

\paragraph{1. Merge Sentiment and Market Data}
The primary merge combined the daily sentiment dataset (\texttt{sentiment\_classified.csv}) with Bitcoin price data (BTC-USD from Yahoo Finance) using matching dates (\texttt{createdAt} $\rightarrow$ \texttt{date}). This alignment linked each tweet’s sentiment score to the corresponding Bitcoin market performance on the same day.

\paragraph{2. Filter for Bitcoin-Relevant Tweets}
After merging, only rows where the \texttt{category} was ``Bitcoin'' were retained. Tweets related to other cryptocurrencies or unrelated topics were removed to focus analysis on Bitcoin-specific sentiment.

\paragraph{3. Clean and Simplify Dataset}
Unnecessary columns (non-relevant flags, usernames, unrelated keyword fields) were dropped to create a compact analysis-ready dataset containing:
\begin{itemize}
    \item \textbf{date / createdAt}
    \item \textbf{sentiment (normalized 0--10)}
    \item \textbf{BTC returns (1D, 7D, 30D)}
    \item \textbf{Tweet volume / count per day}
\end{itemize}

This refined dataset, referred to as \texttt{btc\_sentiment\_daily}, formed the basis for all subsequent analysis and modeling.

\section{Exploratory Data Analysis}

To investigate how Twitter activity relates to Bitcoin market dynamics, we conducted a detailed exploratory data analysis (EDA) incorporating tweet sentiment, tweet volume, and subsequent BTC price behavior. The goal was to visually and statistically inspect whether signals from tweets align with short-horizon Bitcoin market movements.

\subsection{Interactive Tweet-Level Dashboard}

An interactive dashboard was created to facilitate rapid exploration of tweet-specific effects. Users can filter by Tweet ID and view the corresponding:
\begin{itemize}
    \item Tweet content and timestamp
    \item Bitcoin closing price over the following 30 days
    \item Computed next-day, 7-day, and 30-day returns
\end{itemize}

This tool helps identify patterns across individual tweets and their short-term market impact.

\begin{figure}[H]
    \centering
    \includegraphics[width=0.85\textwidth]{dashboard.png}
    \caption{Interactive dashboard showing tweet details and subsequent 30-day BTC price trend.}
    \label{fig:dashboard}
\end{figure}

\subsection{Sentiment vs.\ Next-Day Bitcoin Return}

Daily aggregated sentiment scores were compared against the next-day Bitcoin return to assess whether shifts in sentiment correspond with measurable price reactions. 

A scatter plot with a fitted regression line (Figure~\ref{fig:sentiment_return}) indicates a slight positive relationship, suggesting that more positive tweet sentiment may be associated with marginally higher next-day BTC returns.

\begin{figure}[H]
    \centering
    \includegraphics[width=0.75\textwidth]{sentiment_vs_returns.png}
    \caption{Relationship between tweet sentiment and next-day BTC returns.}
    \label{fig:sentiment_return}
\end{figure}

\subsection{Tweet Volume vs.\ BTC Volatility}

We also examined whether days with higher tweet volume corresponded to greater short-term Bitcoin price volatility. The results, shown in Figure~\ref{fig:volume_volatility}, exhibit a weak positive trend, indicating that increased tweet activity may be loosely associated with heightened BTC volatility.

\begin{figure}[H]
    \centering
    \includegraphics[width=0.75\textwidth]{volatility_vs_volume.png}
    \caption{Tweet volume versus short-term BTC price volatility.}
    \label{fig:volume_volatility}
\end{figure}

\section{Key Research Questions}

\subsection{Descriptive Questions}
\begin{itemize}
    \item How does the overall sentiment about Bitcoin vary over time?
    \item Do peaks in tweet volume correspond to high volatility or large BTC price changes?
\end{itemize}

\subsection{Analytical Questions and Hypotheses}

\begin{table}[h!]
\centering
\renewcommand{\arraystretch}{1.3}
\begin{tabular}{|p{4cm}|p{5cm}|p{5cm}|}
\hline
\textbf{Hypothesis} & \textbf{Null Hypothesis} & \textbf{Test / Approach} \\ \hline
H1: Positive tweet sentiment is associated with positive next-day BTC returns. & 
Tweet sentiment has no relationship with next-day BTC returns. & 
Correlation and regression between sentiment polarity and next-day return. \\ \hline

H2: Negative tweets predict short-term declines in Bitcoin prices (1--7 days). & 
Sentiment does not influence short-term returns. & 
Event study analysis and t-tests of mean returns after negative tweets. \\ \hline

H3: Periods with high tweet activity have higher BTC volatility. & 
Tweet volume is unrelated to BTC volatility. & 
Regression of volatility on tweet volume. \\ \hline

H4: Sentiment effects weaken over time (i.e., 30-day return correlation $<$ 7-day). & 
Sentiment’s effect is consistent over time. & 
Compare $R^2$ or effect sizes across time horizons. \\ \hline
\end{tabular}
\caption{Summary of Analytical Hypotheses and Testing Approaches}
\end{table}

\section{Analytical and Modeling Techniques}

\subsection{Data Transformation}

\begin{itemize}
    \item Convert categorical sentiment to numeric values (e.g., +1, 0, -1).
    \item Aggregate daily tweet metrics:
    \begin{itemize}
        \item Mean sentiment per day
        \item Tweet volume per day
    \end{itemize}
    \item Merge with daily BTC data to create a combined dataset, denoted as \texttt{btc\_sentiment\_daily}.
\end{itemize}

\subsection{Modeling Framework}

\subsubsection{(a) Correlation \& Regression Analysis}

\begin{itemize}
    \item Compute Pearson and Spearman correlations between:
    \begin{itemize}
        \item Daily mean sentiment and next-day returns
        \item Daily mean sentiment and 7-day returns
        \item Daily mean sentiment and 30-day returns
    \end{itemize}

    \item Fit regression models of the form:
   \[
    \mathrm{Return}_{t+h} = \beta_0 + \beta_1\,\mathrm{Sentiment}_t + \beta_2\,\mathrm{Volume}_t + \epsilon_t
\]
where \( h = 1, 7, 30 \) represents different forward-looking time horizons.
\end{itemize}

\section{Deliverables}

\begin{itemize}
    \item \textbf{Statistical Analysis Report:} Hypothesis testing and regression results.
    \item \textbf{Predictive Model:} Forecast or classification model evaluating the predictive power of sentiment.
\end{itemize}

\end{document}
